\documentclass[11pt]{article}
\usepackage[a4paper,margin=27mm]{geometry}
\usepackage[T1]{fontenc}
\usepackage{lmodern,microtype,amsmath,amssymb,amsthm,enumitem,booktabs}
\usepackage[hidelinks]{hyperref}
\newtheorem{theorem}{Theorem}
\newtheorem{proposition}[theorem]{Proposition}
\newtheorem{corollary}[theorem]{Corollary}
\theoremstyle{remark}\newtheorem{remark}[theorem]{Remark}
\newcommand{\Dscr}{\mathcal D}
\newcommand{\M}{\mathcal M}
\newcommand{\C}{\mathcal C}
\newcommand{\supp}{\operatorname{supp}}
\title{Coverage after Neutral Abstention}
\author{}\date{12 October 2026}
\begin{document}\maketitle

\noindent
We adopt one exception: the system covering the vacuous model $V$ may
tie good advice with silence on mixed problems. All other systems retain
strict guidance improvement (G), independence (IB), and closed competitive
selection (S). The exception is limited to $V$ because it is induced by
a change to the background order. Extending the same omission-blind
treatment to every target would give full coverage (see \emph{Omission
Blindness and Four-Rank Guidance}), but that broader weakening is not
adopted. This settles coverage of $V$, but does not settle coverage
of every proper credal set. This note gives a continuous connection between
the exception and ordinary strict rankings, and an exact consistency test
for an ordinal construction. Neither is a full-coverage theorem.
The original account missed two earlier sources. The neutral-abstention
system is already in \texttt{truth\_directed\_losses.tex}, Section 4,
and the constant odds-ratio limit is in \texttt{nonadditive\_losses.tex},
Remark \texttt{twoshun}. The limit theorem below adapts that construction
to the four-rank order and gives an explicit uniform bound; it is not a
new coverage idea. The collaborator has since checked this adaptation and
the consistency argument and reported no error. No human verification
is asserted.

\section{The revised question}

There are finitely many worlds, indexed by $i\le n$, and $\Delta$ is their
probability simplex. Models are measurable advice on binary problems,
identified modulo Gaussian null sets. Count each pair $\{g,-g\}$ once on
$H=\{g:g_1>0\}$, with normalized Gaussian measure $\mu$. Coordinate
hyperplanes can be ignored. A problem is mixed when some coordinates of
$g$ are positive and some negative; otherwise it is unanimous.
At world $i$ the four advice ranks, in order of increasing penalty, are
good, silence, bad, and contradictory.

For a nonempty compact convex credal set $M\subseteq\Delta$, write
\[
 \underline{P_M}(g)=\min_{m\in M}m\cdot g,\qquad
 \Dscr_M=\{g:\underline{P_M}(g)>0\}\cup\{g\ge0:g\ne0\}.
\]
Thus $V=\Dscr_\Delta$, and the ideal at world $i$ is $D_i=\{g:g_i>0\}$,
up to null sets. Smaller $\preceq_i$ means better. Dominance means a weak
improvement at every world and a strict improvement at some world.

Condition (G) preserves pointwise advice improvement, strictly when the
improvement occurs on a non-null set. Condition (IB) preserves comparisons
when advice shared by the two models is replaced on any measurable set
of binary problems. Condition (S) supplies coherent undominated $C_p$,
$p\in\Delta$, such that, for every model $X$, some $i\in\supp p$ has
$C_p\preceq_iX$, and each section $\{p:C_p\preceq_iX\}$ is closed.
Condition (C), when imposed, excludes an incoherent model tied everywhere
with a coherent undominated model.

Coverage (A) concerns a \emph{class} of systems: each coherent target must
be undominated for some system. Our revised question is therefore:
\begin{quote}
For every proper credal set $M\subsetneq\Delta$, does some system satisfying
the standard axioms, including strict (G), (IB), and (S), leave $\Dscr_M$
undominated? Can all these systems be complete, or are incomparabilities
needed for some proper targets?
\end{quote}
The accepted exceptional system adds coverage of $V$ to any positive answer.
The old impossibility result for complete strict rankings does not answer
this revised question: its universally excluded target was $V$.

\section{The exception is a limit of strict rankings}

As the earlier \texttt{twoshun} remark already shows, the exception need
not be viewed as an isolated penalty table. It is the
endpoint of a simple family of strict, independent rankings, whose covered
targets expand towards the whole simplex. This also supplies a useful
positive class in every dimension.

\begin{theorem}\label{thm:limit}
For each $0<\varepsilon\le1$, there are finite real-valued losses on all
models satisfying (G), (IB), (LA), (C), and (S), which cover
\[
 M_\varepsilon
 =\{m\in\Delta:m_i\ge\varepsilon m_j\text{ for all }i,j\}.
\]
These proper credal sets increase to a dense subset of $\Delta$ as
$\varepsilon\downarrow0$, and converge to $\Delta$ in Hausdorff distance.
The losses converge uniformly over models to losses satisfying the agreed
weakening, (IB), (LA), (C), and (S), whose only undominated model is $V$.
\end{theorem}
\begin{proof}
At world $i$, integrate the following penalties:
\[
\begin{array}{c|cccc}
 &\text{good}&\text{silence}&\text{bad}&\text{contradictory}\\\hline
 \text{mixed}&0&\varepsilon|g_i|&(1+\varepsilon)|g_i|
     &2(1+\varepsilon)|g_i|\\
 \text{unanimous}&0&|g_i|&2|g_i|&4|g_i|.
\end{array}
\tag{1}
\]
The integrable bound $4|g_i|$ makes the losses finite. Their strictly
ordered penalties give (G), and cancellation of the common contribution
gives (IB).

For $p\in\Delta$, put
\[
 A=\sum_{g_i>0}p_i g_i,\qquad B=\sum_{g_i<0}p_i(-g_i).
\]
On a mixed problem the expected penalties of $+$, silence, and $-$ are
$(1+\varepsilon)B$, $\varepsilon(A+B)$, and $(1+\varepsilon)A$.
Consequently $+$ uniquely wins when $\varepsilon A>B$, and $-$ when
$\varepsilon B>A$; silence wins between these thresholds. Contradiction
never wins. The tie sets are contained in finitely many proper
hyperplanes and are null, including at boundary priors. On unanimous
problems the good recommendation uniquely wins.

Define
\[
 K_{p,\varepsilon}
 =\left\{\frac{(p_i a_i)_i}{\sum_jp_j a_j}:
                      \varepsilon\le a_i\le1\right\}.
\tag{2}
\]
This set is compact and convex. For convexity, the nonnegative multiples
of the box $\prod_i[\varepsilon p_i,p_i]$ form a convex cone: a sum of
two scaled box points is their total scale times a convex combination
of box points. Intersect that cone with $\Delta$.
The denominator in (2) is positive, and
\[
 \min_{a\in[\varepsilon,1]^n}\sum_i p_i a_i g_i
       =\varepsilon A-B.
\]
Thus the unique minimizing advice is that of the coherent
$C_p=\Dscr_{K_{p,\varepsilon}}$. At the uniform prior, (2) is exactly
$M_\varepsilon$: its points have all coordinate ratios at least
$\varepsilon$, and conversely rescaling by the largest coordinate puts
every such point in $[\varepsilon,1]^n$.

For completeness, the standard Bayes argument verifies the remaining
axioms. Uniqueness of the expected-loss minimizer makes each $C_p$
undominated and makes every model tied with it at all worlds equal to
it modulo null sets. Some supported world must rank $C_p$ no worse than
any given model, by expected-loss optimality. If $p_k\to p$, the unique
minimizing advice at $p$ remains minimizing for all sufficiently large
$k$, at almost every problem. Dominated convergence gives continuity of
each $L_i(C_p)$ and therefore the closed sections in (S). This proves
(LA) and (C) as well.

The sets $M_\varepsilon$ expand as $\varepsilon$ decreases. Every
$q\in\Delta$ is approximated uniformly by the normalization of
$(\max(q_i,\varepsilon))_i$; the $\ell^1$ error is at most $2n\varepsilon$.
This proves the Hausdorff assertion.

Finally put $\varepsilon=0$ in (1). The loss of $V$ is zero at every
world. Every other model either changes unanimous advice, gives a
contradiction, or recommends an option on mixed problems. On a non-null
set such a change has positive penalty at some fixed world. Thus $V$
dominates every other model. The constant selection $C_p=V$ verifies (S),
and (C) and (LA) follow. Good advice and silence tie precisely on mixed
problems; all other adjacent penalty gaps are positive. Moreover
\[
 \sup_X|L_i^\varepsilon(X)-L_i^0(X)|
 \le 2\varepsilon\int_H|g_i|\,d\mu\longrightarrow0.
\]
This establishes uniform convergence and finishes the proof.
\end{proof}

More generally, replace the mixed penalties by
\[
 0,\quad \lambda_i|g_i|,\quad(\lambda_i+\eta_i)|g_i|,
 \quad2(\lambda_i+\eta_i)|g_i|,
 \qquad 0<\lambda_i\le\eta_i<\infty.
\tag{3}
\]
Keep the unanimous penalties in (1). Comparing acceptance with silence
uses $\sum_{g_i>0}p_i\lambda_i g_i-
\sum_{g_i<0}p_i\eta_i(-g_i)$. A positive value also makes acceptance
better than rejection, because $\lambda_i\le\eta_i$.
The same proof gives the normalized boxes
\[
 K_p=\left\{\frac{(p_i a_i)_i}{\sum_jp_j a_j}:
                   \lambda_i\le a_i\le\eta_i\right\}.
\tag{4}
\]
Every normalized box with strictly positive bounds on its supporting
face is covered, including precise models on any face. This is the
global-box construction familiar from concave CPMs, with contradiction
assigned a larger penalty so as to respect the present guidance order.
It is a positive subclass, not a characterization of independent rankings.

\begin{remark}
Uniform convergence of the losses does not preserve strict (G): the
good--silence gap tends to zero. Nor is joint continuity of the selections
in $(p,\varepsilon)$ asserted. In particular $C_{e_i}=D_i$ for
$\varepsilon>0$, while the limiting system may use the constant selection
$V$. These facts explain why the theorem does not contradict the strict
complete-ranking impossibility result.
\end{remark}

\section{A precise ordinal construction problem}

One can specify comparisons directly, without scores, bonuses, or numerical
losses. The following formulation makes the consistency question exact.
It separates choosing closed sections from drawing the consequences of
those choices. It does not assume that suitable choices exist.

Fix a family of coherent models $C_p$, $p\in\Delta$, containing the target.
For every model $X$ choose closed sets $A_i(X)\subseteq\Delta$, subject to
\[
 \text{for every }p\in\Delta,\quad
 p\in A_i(X)\text{ for some }i\in\supp p.
\tag{5}
\]
These are intended to be the \emph{exact} sections
$\{p:C_p\preceq_iX\}$, not merely closed subsets of them.

Start by installing all weak background comparisons and all comparisons
$C_p\preceq_iX$ for $p\in A_i(X)$. Repeatedly apply these four rules:
\begin{enumerate}[leftmargin=*,itemsep=3pt]
\item Install reflexive comparisons.
\item If $X\preceq_iY$ and $Y\preceq_iZ$, install $X\preceq_iZ$.
\item Transfer every comparison under every replacement of shared advice.
\item For each selected $C_p$, if $X\preceq_iC_p$ at all worlds, install
      $C_p\preceq_iX$ at all worlds.
\end{enumerate}
The last rule is precisely what undominatedness requires: a model weakly
better than a protected model must tie with it everywhere.

The following outcomes are forbidden:
\begin{enumerate}[label=(\alph*),leftmargin=*,itemsep=3pt]
\item a reverse of any strict guidance improvement;
\item $C_p\preceq_iX$ with $p\notin A_i(X)$;
\item an incoherent $X$ tied at every world with a selected $C_p$.
\end{enumerate}
For (c), coherence means coherence modulo null sets. Omit (c) if the desired
version of the framework does not impose (C).

\begin{theorem}[Exact consistency test]\label{thm:closure}
For the fixed family and closed sets above, the following are equivalent:
\begin{enumerate}[label=(\roman*),leftmargin=*]
\item There are preorders satisfying (G), (IB), (S), and (C), with that
      selected family and those exact sections.
\item Applying the four rules never produces a forbidden outcome.
\end{enumerate}
When these conditions hold, the generated relations are the least such
relations. They leave every selected model, including the target,
undominated. They need not be complete.

Every failure has a finite derivation from finitely many installed
comparisons. Applying all rules in parallel at stages $0,1,2,\ldots$ and
taking the union gives the full closure; no later ordinal stages add
anything once the family and sections are fixed.
\end{theorem}
\begin{proof}
Every proposed solution contains the initial comparisons and is closed
under the four rules. For the fourth rule, otherwise $X$ would dominate
the selected $C_p$. Such a solution cannot have any forbidden outcome.
Thus (i) implies (ii), and every generated comparison is necessary.

Conversely, take the generated relations. The first two rules give
preorders. The third gives (IB) in both directions, since the inverse
replacement is also allowed. The initial comparisons and prohibition (a)
give weak and strict (G). The fourth rule makes each $C_p$ undominated.
Initial installation and prohibition (b) make its sections exactly the
closed sets $A_i(X)$; (5) gives competitiveness. This proves (S) and (LA).

To verify (C) for any coherent undominated model, not just a selected
one, use the simplex consequence of (S): every model has a selected weak
dominator. Indeed the closed sections for that model cover each face by
its own indices, so the simplex lemma gives a point in their intersection.
An undominated model must tie everywhere with this selected dominator.
A model tied everywhere with it then ties with the same selected model,
and prohibition (c) makes it coherent. The tie also makes it undominated.

Finally, each rule has finitely many premises, since $n$ is finite.
Any premises present in the union of the finite stages already occur
together at one finite stage, so their conclusion occurs at the next.
The union is therefore closed under all four rules. Induction on the
stage gives a finite proof tree for each comparison, and each forbidden
outcome involves only finitely many comparisons. Leastness follows from
the first paragraph.
\end{proof}

There is no loss in requiring the target to occur in the family.
In any successful system, the simplex argument makes an undominated target
tie everywhere with some $C_r$. Replace that $C_r$ by the target.
Transitivity makes all comparisons with every $X$ unchanged, so none of
the sections, competitiveness conditions, or undominatedness properties
changes.

\paragraph{What this says about transfinite induction.}
The order-theoretic closure is not the unresolved part. Once closed
sections are chosen consistently, the theorem constructs the required
partial rankings by ordinary iteration. What remains is to choose the
family and its closed sections so that no finite contradiction is forced.
A transfinite construction could still make those choices one at a time,
but it needs an extension lemma preserving their eventual closedness.
An increasing union of closed sections need not be closed, and adding
missing limit comparisons may trigger the forbidden outcomes above.
The theorem neither supplies that extension lemma nor makes checking all
finite derivations an effective finite algorithm.

\section{What is resolved and what remains open}

\begin{itemize}[leftmargin=*,itemsep=5pt]
\item \textbf{Resolved: coverage of $V$ under the agreed exception.}
The limiting system above is complete and independent. At world $i$,
$D_i$ ties with $V$. Good--silence indifference on mixed problems holds
in every common context, as intended.
\item \textbf{Resolved: two worlds.} Version 3 already covers every
credal interval under strict (G), (IB), and (S). The exception can be
adjoined, or used for $V$, without losing any of that coverage.
\item \textbf{Resolved: a positive class in every dimension.}
Precise models and the normalized boxes in (4) have complete independent
rankings satisfying all the standard requirements. Theorem~\ref{thm:limit}
also covers a nested family approaching $V$.
\item \textbf{Earlier pricing restrictions.} Section 2 of
\texttt{truth\_directed\_losses.tex} excludes targets with incomparable
boundary supports under its exchange-rate and domain assumptions,
including suitable nonadditive losses, already in two worlds. The CPM
boundary-support theorem in v19b excludes different boundary supports,
including nested supports. These are additional restrictions on particular
methods; the exchange-rate assumption is not implied by the present
ordinal axioms.
\item \textbf{Still impossible: full revised coverage by additive losses
in $n\ge4$.} The v3 obstruction excludes every credal target containing
three vertices. Some are proper, for example
$\operatorname{conv}\{e_1,e_2,e_3\}$ in four worlds. The finitely additive
one-sided-extension obstruction in \emph{Closed Choice Beyond Ordinary
Integrals} gives the same restriction under its stated hypotheses.
The $V$ exception does not apply to these proper targets.
\item \textbf{Open: coverage of every proper target in $n\ge3$.}
Neither the existing results nor this note establishes it for complete
rankings or for incomplete ones. The complete-ranking clique theorem
narrows targets containing three vertices, but does not exclude all
proper such targets. For other targets, the failure of a particular
CPM or box construction is not an impossibility theorem for ordinal
rankings.
\item \textbf{Open: the consistent choice required by
Theorem~\ref{thm:closure}.} This is a precise remaining task for an
ordinal construction. Complete rankings would require additional choices
of orientations for every currently incomparable pair; the theorem does
not justify making those choices.
\end{itemize}

In particular, incomparabilities are \emph{not currently known to be
necessary} for the revised coverage demand. Their necessity for the old
demand followed from excluding $V$ under strict (G), and that argument
no longer applies. Nor have incomparabilities been proved sufficient.

\medskip\noindent
\textbf{Guide to the existing proofs.} The companion file
\emph{START HERE -- Coverage Status.md} records the current questions,
proof locations, dependencies, and historical notes that should no longer
be read as statements of the current open problem. It leaves the original
papers in place so existing references continue to work.
\end{document}
